\section{Experimental Evaluation}
\label{sec:evaluation}

This section evaluates MPTA-Repair on both benchmark and industrial
case studies. The evaluation follows the clock-bound repair setting
defined in Section~\ref{sec:approach}, where only numerical bounds in
guards and invariants are repairable. Comparing MPTA-Repair with an
iterative baseline extended from the TARTAR paradigm
\cite{leue2019clock} demonstrates that multi-property regression and
multi-counterexample analysis are essential to prevent secondary
property violations.

\subsection{Evaluation Strategy and Datasets}

We evaluate MPTA-Repair on two datasets, summarized in
Table~\ref{tab:dataset-statistics}. The benchmark dataset includes nine
standard models from UPPAAL \cite{behrmann2004tutorial} and the
XtaBenchmarkSuite \cite{farkas2018benchmarks}, such as Elevator and
FDDI. The industrial dataset contains five closely coupled industrial
case studies: Gear Control System for automotive transmission
\cite{lindahl1998gear}, SAE RoR for autonomous-driving decisions
\cite{alves2021double}, Train Controller Alarm for railway emergency
timing \cite{gonzalez2024uppaal}, Train Gate Controller for
shared-resource concurrency \cite{alur1994theory}, and Pacemaker for
medical device pacing \cite{jiang2010heart}. Compared with the
benchmark models, these case studies involve more tightly coupled
timing constraints across multiple properties.

\begin{table}[t]
\caption{Statistics of the benchmark and industrial datasets.}
\label{tab:dataset-statistics}
\centering
\scriptsize
\setlength{\tabcolsep}{2pt}
\begin{tabular*}{\linewidth}{@{\extracolsep{\fill}}lccccc@{}}
\toprule
Dataset & Original models & Mutated models & Avg. properties/model &
Avg. loc. & Avg. trans. \\
\midrule
Benchmark & 9 & 54 & 4.1 & 30.56 & 52.67 \\
Industrial & 5 & 105 & 7.3 & 57.80 & 91.80 \\
\bottomrule
\end{tabular*}
\end{table}


Since real-world timed-automata models with labeled clock-bound faults
are difficult to obtain, we synthesize faulty instances using mutation
\cite{jia2011mutation}. We retain only mutants that are syntactically
valid, violate at least one timing safety property, and preserve the
untimed behavior of the original model.

We compare MPTA-Repair against an adapted iterative TARTAR-style
baseline \cite{leue2019clock}. Because the TARTAR method was originally
designed for single-property repair and lacks joint multi-counterexample
relation analysis, this comparison allows us to examine whether using a
joint counterexample pool improves repair robustness over trace-local
repair generation. The comparison is designed to evaluate the effect of
jointly analyzing properties, counterexamples, and repair variables.

\subsection{Experimental Results}

The experimental evaluation demonstrates that MPTA-Repair outperforms
the iterative baseline, particularly on complex industrial models. As
shown in Table~\ref{tab:overall-results}, the baseline achieves a 72\%
success rate on the benchmark dataset, whereas MPTA-Repair reaches
93\%. This performance gap widens on the industrial dataset, where the
baseline repairs only 20\% of the cases compared to 76\% for
MPTA-Repair. These results indicate that our approach remains effective
on standard timed-automata benchmarks while demonstrating robustness in
complex industrial systems where properties and traces interact closely.
The average repair time should be interpreted together with the success
counts. On the industrial dataset, the iterative baseline has a lower
average time because many cases fail quickly, not because it repairs
industrial models more efficiently.

\begin{table}[t]
\caption{Repair success rates across benchmark and industrial datasets.}
\label{tab:overall-results}
\centering
\scriptsize
\setlength{\tabcolsep}{4pt}
\begin{tabular*}{\linewidth}{@{\extracolsep{\fill}}llc@{}}
\toprule
Dataset & Method & Success rate \\
\midrule
Benchmark & Iterative baseline & 72\% \\
Benchmark & MPTA-Repair & 93\% \\
Industrial & Iterative baseline & 20\% \\
Industrial & MPTA-Repair & 76\% \\
\bottomrule
\end{tabular*}
\end{table}


The baseline is competitive on simpler benchmark models because timing
errors are frequently concentrated on isolated guards or invariants. In
these localized scenarios, a single diagnostic trace often provides
sufficient information to infer a correct repair. However, even within
standard benchmarks featuring shared multi-channel structures, such as
the FDDI protocol, MPTA-Repair is more robust because it prevents
single-trace overfitting. This advantage becomes more pronounced in
complex industrial models that feature interacting concurrent
components, urgent or broadcast synchronization channels, observer
templates, and array-based constructs. In these settings, localized
repairs by the baseline frequently fail subsequent global validation
checks. Our empirical analysis shows that the baseline lacks explicit
optimization to minimize edit magnitude, often generating extreme
candidate bounds that violate admissibility. Additionally, the baseline
struggles with regression verification over three to nine coupled
properties because it does not analyze their corresponding
counterexamples jointly.

MPTA-Repair achieves higher success rates due to two primary design
advantages. First, multi-counterexample accumulation reduces overfitting
to a single timed diagnostic trace, which typically exposes only one
timing realization of an underlying fault. By accumulating and analyzing
diagnostic traces from multiple operational modes, the framework
synthesizes minimal joint edit sets that address the timing fault
comprehensively. Second, relation-guided optimization complements
precise MaxSMT objectives to improve candidate quality. By identifying
diagnostic traces that share path fragments or target repair variables,
the solver focuses on modifying numerical clock bounds close to the
observed violations while penalizing large edit magnitudes. Successful
repairs therefore typically require modifying only one or two critical
clock bounds, which facilitates subsequent manual inspection and
verification.

\subsection{Discussion and Practical Lessons}

The evaluation highlights several insights for applying automated repair
to industrial timed automata. First, full-property regression testing is
necessary. Because industrial system properties are interdependent, a
localized fix designed solely to satisfy a liveness-like response
monitor can inadvertently violate a separate safety monitor that
constrains maximum waiting times. MPTA-Repair avoids these behavioral
regressions by treating overall property restoration as a global
objective. Second, admissibility is critical to measuring repair
success. Without ensuring structural and functional equivalence, a
solver might satisfy a specific timing property by disabling the
intended functional behavior, such as preventing a critical system
action entirely. The strict admissibility check ensures that property
satisfaction is not achieved at the expense of untimed functional
behavior.

Our failure analysis reveals several limitations of the current repair
space. Computational timeouts and unsupported model syntax account for
only a portion of the unsuccessful repair attempts. In tightly coupled
domains such as cardiac pacing, certain complex errors require
simultaneous, consistent adjustments across both guards and invariants,
which significantly expands the search space. Additionally, incomplete
diagnostic trace coverage from the verification engine can occasionally
lead to overfitting. Certain complex cases lacked any feasible candidate
within the restricted clock-bound parameter space. This limitation
suggests that repairing these models requires broader structural
modifications, such as altering synchronization labels, adjusting clock
resets, or modifying discrete data updates.
